% Core packages for Beamer slides, Japanese hyperlinks, figures, tables, and diagrams.
\usepackage{bxdpx-beamer}
\usepackage{pxjahyper}
\usepackage{graphicx}
\usepackage{booktabs}
\usepackage{array}
\usepackage{xcolor}
\usepackage{tikz}
\usepackage[normalem]{ulem}
\usepackage{listings}
\usepackage{csquotes}
\usepackage{ifthen}

\lstset{
  basicstyle=\ttfamily\tiny,
  breaklines=true,
  frame=single,
  backgroundcolor=\color{black!5},
  keywordstyle=\color{blue},
  commentstyle=\color{green!50!black},
  stringstyle=\color{red}
}

% Define new command for quoting text with double quotes.
\newcommand{\qbox}[1]{\mbox{`#1'}}
\newcommand{\qqbox}[1]{\mbox{``#1''}}

% original \href  -> \oldhref
\let\oldhref\href
% renewcommand \href 
\renewcommand{\href}[2]{%
  \oldhref{#1}{\textcolor{blue}{\underline{#2}}}%
}
% Make arXiv links blue and underlined
\newcommand{\arxiv}[1]{\href{https://arxiv.org/abs/#1}{arXiv:#1}}


% Footnote setup: dagger-style labels and compact formatting for Beamer slides.
\newcounter{presentationfootnote}
\renewcommand{\thefootnote}{\mbox{$\dagger$}\arabic{footnote}}
\renewcommand{\thempfootnote}{\mbox{$\dagger$}\arabic{mpfootnote}}
\renewcommand{\thepresentationfootnote}{\mbox{$\dagger$}\arabic{presentationfootnote}}
\setbeamertemplate{footnote mark}{\textsuperscript{\insertfootnotemark}}

% Format the footnote body with a compact label column and left-aligned text.
\setbeamertemplate{footnote}{%
  \parindent 0pt%
  \noindent%
  \raggedright%
  \hbox to 1.3em{\hfil\insertfootnotemark}%
  \insertfootnotetext\par%
}

% Redefine Beamer footnotes so dagger-style labels also work with overlay syntax.
\makeatletter
\renewcommand<>{\footnote}[1][]{%
  \let\beamer@footnotetext=\@footnotetext%
  \let\beamer@mpfn=\@mpfn%
  \let\beamer@thempfn=\thempfn%
  \let\beamer@kvorig=\KV@errx%
  \let\beamer@xkvorig=\XKV@err
  \def\beamer@footarg{}%
  \def\KV@errx##1{\edef\beamer@footarg{\@tempa}}%
  \def\XKV@err##1{\edef\beamer@footarg{\XKV@tkey}}%
  \ifbeamer@framefootnotes
    \setkeys{beamerfootnote}{frame,#1}%
  \else
    \setkeys{beamerfootnote}{#1}%
  \fi
  \let\KV@errx=\beamer@kvorig%
  \let\XKV@errx=\beamer@xkvorig
  \ifx\beamer@footarg\@empty%
    \def\beamer@next{\stepcounter{presentationfootnote}%
      \protected@xdef\@thefnmark{\thepresentationfootnote}%
      \@footnotemark\beamer@footnotetext#2}%
  \else%
    \def\beamer@next{%
      \begingroup
        \unrestored@protected@xdef\@thefnmark{\mbox{$\dagger$}\beamer@footarg}%
      \endgroup
      \@footnotemark\beamer@footnotetext#2}%
  \fi%
  \beamer@next}
\makeatother

% Base theme: keep the deck simple and hide default navigation icons.
\usetheme{default}
\usecolortheme{default}
\setbeamertemplate{navigation symbols}{}

% Footer: page number shown slightly above the bottom-right corner.
\setbeamertemplate{footline}{
  \leavevmode
  \hbox to \paperwidth{%
    \hfill
    \raisebox{0.95em}[0pt][0pt]{%
      \usebeamerfont{page number in head/foot}%
      \usebeamercolor[fg]{page number in head/foot}%
      \color{white}%
      \insertframenumber\,/\,\inserttotalframenumber\hspace{0.8em}%
    }%
  }
}

% Frame title style.
\setbeamercolor{frametitle}{fg=darkgray}
\setbeamerfont{frametitle}{series=\bfseries,size=\Large}

% Caption formatting.
\renewcommand{\figurename}{Figure}
\setbeamercolor{caption name}{fg=black}
\setbeamertemplate{caption}[numbered]
\setbeamertemplate{caption label separator}{: }
\setlength{\abovecaptionskip}{-3pt}
\setlength{\belowcaptionskip}{0pt}

% Clickable paper-source link helper used near the frame title.
\newcommand{\sourcelink}[2]{\oldhref{#2}{\textcolor{blue!65!black}{\uline{\texttt{#1}}}}}

% Custom frame title with separator rule and top-right paper reference.
\setbeamertemplate{frametitle}{
  \begin{flushleft}
    \vspace{0.01em}
    \hspace{0.1cm}\insertframetitle\par
    \vspace{0.07em}
    \hrule width 0.89\linewidth
    \vspace{-0.2em}
  \end{flushleft}
  \begin{tikzpicture}[remember picture,overlay]
    \node[anchor=north east, inner sep=0pt] at ([xshift=-0.55cm,yshift=-1.13cm]current page.north east)
      {\scriptsize\normalfont\ifcsname papersource\endcsname\papersource\fi};
  \end{tikzpicture}
  \vspace{0.55em}
}

% Block appearance for normal, example, and alert emphasis.
\setbeamertemplate{blocks}[rounded][shadow=true]
\setbeamercolor{block title}{fg=white,bg=blue!55!black}
\setbeamercolor{block body}{fg=black,bg=blue!6}
\setbeamercolor{block title example}{fg=white,bg=green!39!black}
\setbeamercolor{block body example}{fg=black,bg=green!10}
\setbeamercolor{block title alerted}{fg=white,bg=red!70!black}
\setbeamercolor{block body alerted}{fg=black,bg=orange!10}
\setbeamercolor{block title code}{fg=white,bg=gray!60!black}
\setbeamercolor{block body code}{fg=black,bg=gray!10}

% Define confidential. 
\newcommand{\confidentialstamp}[1][en]{%
  \ifthenelse{\equal{#1}{ja}}{%
    \def\confidentialtext{社外秘}%
  }{%
    %\def\confidentialtext{CONFIDENTIAL}%
    \def\confidentialtext{\resizebox{1.1cm}{0.2cm}{CONFIDENTIAL}}%
  }%
  \begin{tikzpicture}[remember picture,overlay]
    \node[
      anchor=north west,
      draw=red!75!black,
      text=red!75!black,
      line width=1.2pt,
      rounded corners=1pt,
      inner xsep=1pt,
      inner ysep=1pt,
      font=\bfseries\scriptsize,
      rotate=+13, 
      opacity=0.9
    ] at ([xshift=0.5mm,yshift=-2.9mm]current page.north west)
    {\confidentialtext};
  \end{tikzpicture}%
}
